\subsection{Jointly Distributed Random Variables}
\hrulefill

\paragraph{Definition}
Two random variables $X$ and $Y$ are said to be \textbf{jointly distributed} if their outcomes are determined 
by the same random experiment. The \textbf{joint probability distribution} of $X$ and $Y$ gives the probabilities 
associated with all possible pairs of outcomes of $X$ and $Y$. The pmf or pdf of the joint distribution is denoted by
\[
P(x,y) = P(X=x, Y=y)
\]

\paragraph{Example}
Suppose you have a hat with 3 balls labeled 1, 2, and 3. You randomly select one ball, record its number as $X$,
then flip a coin that many times and record the number of heads as $Y$. The possible outcomes of this experiment are:\\

a. Find $P(1,0)$:
\[
P(1,0) = P(X=1, Y=0) = \frac{1}{3} \cdot \frac{1}{2} = \frac{1}{6}
\]

b. Find $P(1,1)$:
\[
P(1,1) = P(X=1, Y=1) = \frac{1}{3} \cdot \frac{1}{2} = \frac{1}{6}
\]

c. Find $P(2,0)$:
\[
P(2,0) = P(X=2, Y=0) = \frac{1}{3} \cdot \frac{1}{4} = \frac{1}{12}
\]

d. Find $P(2,1)$:
\[
P(2,1) = P(X=2, Y=1) = \frac{1}{3} \cdot \frac{1}{2} = \frac{1}{6}
\]

e. Find $P(2,2)$:
\[
P(2,2) = P(X=2, Y=2) = \frac{1}{3} \cdot \frac{1}{4} = \frac{1}{12}
\]

f. Find $P(3,0)$:
\[
P(3,0) = P(X=3, Y=0) = \frac{1}{3} \cdot \frac{1}{8} = \frac{1}{24}
\]  

g. Find $P(3,1)$:
\[
P(3,1) = P(X=3, Y=1) = \frac{1}{3} \cdot \frac{3}{8} = \frac{1}{8}
\]

h. Find $P(3,2)$:
\[
P(3,2) = P(X=3, Y=2) = \frac{1}{3} \cdot \frac{3}{8} = \frac{1}{8}
\]

i. Find $P(3,3)$:
\[
P(3,3) = P(X=3, Y=3) = \frac{1}{3} \cdot \frac{1}{8} = \frac{1}{24}
\]

\subsubsection*{Marginal PMFs}
The \textbf{marginal pmf} of a random variable is the pmf of that variable without reference to the other variable.
We denote the marginal pmfs of $X$ and $Y$ as $P_X(x)$ and $P_Y(y)$, respectively. They can be found by summing
over the joint pmf:
\[
P_X(x) = \sum_y P(x,y)
\]
\[
P_Y(y) = \sum_x P(x,y)
\]

\paragraph{Example (continued)}
Find the marginal pmfs of $X$ and $Y$ from the previous example.
\begin{align*}
    P_X(1) &= P(1,0) + P(1,1) = \frac{1}{6} + \frac{1}{6} = \frac{1}{3} \\
    P_X(2) &= P(2,0) + P(2,1) + P(2,2) = \frac{1}{12} + \frac{1}{6} + \frac{1}{12} = \frac{1}{3} \\
    P_X(3) &= P(3,0) + P(3,1) + P(3,2) + P(3,3) = \frac{1}{24} + \frac{1}{8} + \frac{1}{8} + \frac{1}{24} = \frac{1}{3}\\
    P_Y(0) &= P(1,0) + P(2,0) + P(3,0) = \frac{1}{6} + \frac{1}{12} + \frac{1}{24} = \frac{7}{24} \\
    P_Y(1) &= P(1,1) + P(2,1) + P(3,1) = \frac{1}{6} + \frac{1}{6} + \frac{1}{8} = \frac{11}{24} \\
    P_Y(2) &= P(2,2) + P(3,2) = \frac{1}{12} + \frac{1}{8} = \frac{5}{24} \\
    P_Y(3) &= P(3,3) = \frac{1}{24}
\end{align*}

\paragraph*{Trick}
For any $X, Y$,
\[
P(X=Y) + P(Y > X) + P(X > Y) = 1
\]

\hrulefill

\subsubsection*{Independence}
Two discrete random variables $X$ and $Y$ are said to be \textbf{independent} if for all values of $x$ and $y$,
\[
\forall (x,y) \in \mathbb{R}, \quad P(X=x, Y=y) = P_X(x) \cdot P_Y(y)
\]

\pagebreak

\subsubsection*{Conditional Distributions}
\paragraph*{Definition} The \textbf{conditional pmf} of $Y$ given $X=x$ is defined as
\[P_{Y|X}(y|x) = P(Y=y | X=x) = \frac{P(X=x, Y=y)}{P_X(x)}\]
Similarly, the conditional pmf of $X$ given $Y=y$ is defined as
\[P_{X|Y}(x|y) = P(X=x | Y=y) = \frac{P(X=x, Y=y)}{P_Y(y)}\]

\paragraph*{Example (continued)}:

a. Find $P_{X|Y}(3|y=2)$:
\[
P_{X|Y}(3|2) = \frac{P(3,2)}{P_Y(2)} = \frac{\frac{1}{8}}{\frac{5}{24}} = \frac{3}{5}
\]

b. Find $P_{Y|X}(3|x=2)$:
\[
P_{Y|X}(3|2) = \frac{P(2,3)}{P_X(2)} = \frac{0}{\frac{1}{3}} = 0
\]

c. Find $P_{Y|X}(1|x=3)$:
\[
P_{Y|X}(1|3) = \frac{P(3,1)}{P_X(3)} = \frac{\frac{1}{8}}{\frac{1}{3}} = \frac{3}{8}
\]

d. Find conditional pmf of $X$ given $Y=1$:
\begin{align*}
    P_{X|Y}(1|1) &= \frac{P(1,1)}{P_Y(1)} = \frac{\frac{1}{6}}{\frac{11}{24}} = \frac{4}{11} \\
    P_{X|Y}(2|1) &= \frac{P(2,1)}{P_Y(1)} = \frac{\frac{1}{6}}{\frac{11}{24}} = \frac{4}{11} \\
    P_{X|Y}(3|1) &= \frac{P(3,1)}{P_Y(1)} = \frac{\frac{1}{8}}{\frac{11}{24}} = \frac{3}{11}
\end{align*}